\documentclass[12pt,a4paper]{article}

\usepackage{iftex}
\ifPDFTeX
  \usepackage[T2A]{fontenc}
  \usepackage[utf8]{inputenc}
  \usepackage[russian,english]{babel}
\else
  \usepackage{fontspec}
  \IfFontExistsTF{Times New Roman}
    {\setmainfont{Times New Roman}}
    {\setmainfont{TeX Gyre Termes}}
  \usepackage{polyglossia}
  \setdefaultlanguage{russian}
  \setotherlanguage{english}
\fi

\usepackage{geometry}
\geometry{left=25mm,right=20mm,top=20mm,bottom=20mm}
\usepackage{booktabs}
\usepackage{hyperref}
\usepackage{xurl}

\title{MergeMind: генерация и оценка review-комментариев для поддержки coding-agent в CI-цикле}
\author{<ФИО студента> \and <ФИО научного руководителя>}
\date{Черновик, 2026}

\begin{document}

\maketitle

\begin{abstract}
В статье рассматривается подход к автоматизированной генерации review-комментариев
по diff и контексту репозитория. Предлагаемый прототип MergeMind использует
цепочку generator--reranker--rewriter для получения малого числа конкретных
комментариев, после чего комментарии оцениваются как в задаче review, так и
через downstream-поведение coding-agent в SWE-CI. Первичные эксперименты
показывают, что MergeMind пока не доказывает сокращение числа итераций, но может
улучшать официальную метрику SWE-CI EvoScore на отдельных задачах и снижать
финальный test gap.
\end{abstract}

\section{Введение}

Code review остается одним из основных механизмов контроля качества
программного обеспечения. Современные языковые модели позволяют генерировать
естественно-языковые замечания по diff, однако практическая ценность такого
инструмента определяется не количеством комментариев, а их полезностью,
обоснованностью и способностью привести к исправлению.

Цель работы --- исследовать, может ли автоматический review-комментарий
использоваться как дополнительный контекст для coding-agent внутри CI-подобного
цикла решения задачи.

\section{Обзор литературы}

Работы по automated code review рассматривают несколько связанных задач:
оценку качества изменения, генерацию комментария и исправление кода по
комментарию. CodeReviewer формулирует эту область как отдельную задачу
предобучения для работы с diff и review comments. CodeReviewQA смещает фокус на
понимание логики review-комментария, а не только на текстовое сходство.

Для оценки естественно-языковых выводов используются как классические метрики
сходства, так и LLM judge с рубриками groundedness и usefulness. Для проверки
практического эффекта review-комментариев полезны CI-like benchmark, где
комментарий можно передать агенту-разработчику и проверить изменение поведения
по тестам.

\section{Постановка задачи}

Дан diff изменения, контекст репозитория и, в SWE-CI сценарии, требование
текущей итерации в виде \texttt{requirement.xml}. Необходимо сгенерировать
небольшое число комментариев, которые:

\begin{itemize}
  \item основаны только на доступном diff и видимом контексте;
  \item указывают на конкретный риск корректности;
  \item пригодны для передачи агенту-разработчику как контекст исправления;
  \item не используют target commit или скрытую oracle-информацию.
\end{itemize}

\section{Предложенный подход}

MergeMind строится как цепочка из нескольких агентов:

\begin{itemize}
  \item generator формирует candidate comments по diff и контексту;
  \item reranker отбирает наиболее полезные и grounded comments;
  \item rewriter переводит выбранные замечания в компактный repair contract;
  \item optional judge используется только для оценки, а не для подсказки SWE-CI.
\end{itemize}

В assisted SWE-CI режиме комментарий передается внутрь цикла решения. Сначала
SWE-CI programmer делает первичную правку, затем MergeMind анализирует
before/after diff и формирует review. После этого review копируется в контейнер
как \texttt{/app/mergemind\_review.md}; programmer выполняет revision pass, и
тесты запускаются уже после revision.

\section{Экспериментальная схема}

Сравниваются два режима:

\begin{itemize}
  \item baseline --- обычный SWE-CI без MergeMind, но с той же локальной моделью;
  \item assisted --- SWE-CI с дополнительным MergeMind review-context между
    programmer patch и запуском тестов.
\end{itemize}

Для оценки используются:

\begin{itemize}
  \item \texttt{actual\_iterations};
  \item official SWE-CI EvoScore;
  \item final gap и best gap;
  \item overlap наборов падающих тестов;
  \item token usage и число LLM-вызовов.
\end{itemize}

\section{Начальные результаты}

\begin{center}
\begin{tabular}{lrrrr}
  \toprule
  Run & EvoScore & Final gap & Comments & Revisions \\
  \midrule
  Baseline \texttt{inline-snapshot} & 0.5045 & 12 & 0 & 0 \\
  Contract \texttt{inline-snapshot} & 0.9167 & 13 & 9 & 3 \\
  Triage \texttt{inline-snapshot} & 0.8333 & 12 & 2 & 2 \\
  Baseline \texttt{cle-b/httpdbg} & -0.5439 & 11 & 0 & 0 \\
  Triage-control \texttt{cle-b/httpdbg} & 0.0000 & 5 & 1 & 1 \\
  Caveman top1 \texttt{cle-b/httpdbg} & -0.3333 & 5 & 2 & 2 \\
  \bottomrule
\end{tabular}
\end{center}

На задаче \texttt{inline-snapshot} MergeMind-профили дали более высокий
EvoScore, чем baseline без MergeMind. На задаче \texttt{cle-b/httpdbg}
положительный EvoScore пока не достигнут, но assisted-варианты снизили final
gap с 11 до 5 и не добавили новых финальных падений в лучшем caveman-профиле.

\section{Обсуждение}

Промежуточные результаты показывают, что эффект зависит от типа комментария.
Более строгие комментарии, ориентированные на один конкретный риск
корректности, дают лучший баланс качества и стоимости, чем более широкие
много-комментарные варианты. При этом сокращение числа итераций пока не
доказано: завершенные прогоны обычно доходили до заданного лимита итераций.

\section{Заключение и дальнейшая работа}

Текущий прототип подтверждает реализуемость in-loop review-context для SWE-CI.
Следующий этап --- расширить число проектов и итераций, отдельно проанализировать
изменение \texttt{requirement.xml} между итерациями и подготовить таблицу
результатов на непересекающемся наборе проектов.

\begin{thebibliography}{9}

\bibitem{codereviewer}
Li Z., Lu S., Guo D. et al. Automating Code Review Activities by Large-Scale
Pre-training. arXiv:2203.09095, 2022.

\bibitem{towards-acr}
Tufano M., Watson C., Bavota G., Poshyvanyk D. Towards Automating Code Review
Activities. arXiv:2101.02518, 2021.

\bibitem{codereviewqa}
CodeReviewQA: The Code Review Comprehension Assessment for Large Language
Models. Findings of ACL, 2025.

\bibitem{geval}
Liu Y. et al. G-Eval: NLG Evaluation using GPT-4 with Better Human Alignment.
arXiv:2303.16634, 2023.

\bibitem{sweci}
SWE-CI: Evaluating Agent Capabilities in Maintaining Codebases via Continuous
Integration. arXiv:2603.03823, 2026.

\end{thebibliography}

\end{document}
